\chapter{Hall pre-check: il livello di benzina prima di partire}
\label{ch:metodo_hall}

\epigraph{``A necessary and sufficient condition that all the
sets in the family $S$ have a system of distinct
representatives is that for every $k$, the union of any
$k$ sets of $S$ contains at least $k$ elements.''}{Philip
Hall, 1935~\parencite{hall1935}}

\lettrine[lines=3]{P}{rima} di accendere la macchina per un
viaggio lungo, qualsiasi automobilista controlla se c'\`e
abbastanza benzina. Sembra ovvio. Eppure quando si lancia un
solver di timetabling \`e una pratica che spesso si salta:
si parte, si aspetta dieci minuti, e si scopre che il
problema era \emph{infeasibile dall'inizio}. Tutti i tentativi
del solver erano destinati a fallire perch\'e i dati di input
non potevano produrre nessun orario valido.

Il pre-check di Hall \`e quel controllo del livello di
benzina: in pochi millisecondi, prima di lanciare CP-SAT,
verifica se le condizioni necessarie alla fattibilit\`a sono
soddisfatte. Se Hall fallisce, $\pi$Tantum lo dice
all'utente (e indica \emph{quale} sotto-insieme di entit\`a \`e
il problema), evitando lunghe attese inutili.

\section{Il teorema dei matrimoni}

\begin{storiabox}
Nel 1935 Philip Hall pubblic\`o sul \emph{Journal of the
London Math. Society}~\parencite{hall1935} un teorema
deliziosamente semplice. Considera un gruppo di $n$ ragazze
e $m$ ragazzi; ogni ragazza ha una lista di ragazzi che
sarebbe disposta a sposare. Quando \`e possibile sposare
\emph{tutte} le ragazze in modo che ognuna abbia un suo
ragazzo distinto?

La risposta di Hall: \emph{quando per ogni sotto-insieme di
$k$ ragazze, la lista combinata di candidati contiene
almeno $k$ ragazzi}.

\`E condizione necessaria e sufficiente. Da allora il
risultato si chiama \emph{teorema dei matrimoni di Hall}, o
\emph{Hall's marriage theorem}\index{Hall!teorema dei
matrimoni}.
\end{storiabox}

\section{Dal teorema all'orario}

L'analogia con il timetabling \`e diretta. Considera la
versione semplificata in cui devi assegnare ad ogni
\emph{lezione} (combinazione classe-materia-ora) un
\emph{docente} compatibile (con la classe di concorso giusta
e la disponibilit\`a sufficiente).

\begin{defbox}
Costruiamo un grafo bipartito\index{matching!bipartito}:
\begin{itemize}
\item A sinistra: le \emph{ore-cattedra} richieste (es.
  ``2 ore di matematica per la 3A'').
\item A destra: i \emph{docenti} disponibili.
\item Un arco $\{l, d\}$ se il docente $d$ \`e compatibile
  con la lezione $l$ (giusta classe di concorso, ore
  residue sufficienti, etc.).
\end{itemize}
Una soluzione di Phase~A esiste se e solo se questo grafo
ha un \emph{matching che copre tutte le ore-cattedra}.
\end{defbox}

\begin{ideabox}
Per il teorema di Hall, il matching esiste \textbf{se e
solo se} per ogni sotto-insieme $S$ di ore-cattedra, la
lista combinata dei docenti compatibili con almeno una di
esse contiene un \emph{numero di ore disponibili almeno
pari} a $|S|$.

In linguaggio operativo: \emph{ogni gruppo di lezioni che si
contendono un pool di docenti compatibili deve avere
abbastanza ore complessive nel pool}.
\end{ideabox}

\section{Tre controlli pratici}

Verificare la condizione di Hall \emph{esattamente} richiede
di esaminare \emph{tutti} i sottoinsiemi di lezioni:
$2^n$ con $n$ ore-cattedra, infattibile per $n$ grande.
$\pi$Tantum implementa allora tre controlli rapidi che,
combinati, intercettano la stragrande maggioranza degli
infeasibility che si vedono in pratica.

\subsection{Controllo 1: per (classe, materia)}

\begin{defbox}
Per ogni coppia $(c, m)$ con $h_{cm}$ ore richieste,
verifica che la somma delle ore disponibili dei docenti
abilitati a $m$ \emph{nella classe di concorso pertinente
per} $c$ sia $\geq h_{cm}$.
\end{defbox}

Tipico caso di fallimento: ``servono 5 ore di matematica nel
liceo classico ma i due docenti di A027 hanno gi\`a
saturato le loro ore in altri istituti''. $\pi$Tantum lo
intercetta immediatamente e mostra il deficit.

\subsection{Controllo 2: subject supply vs demand}

\begin{defbox}
Per ogni materia $m$, la \emph{domanda} \`e la somma di
$h_{cm}$ su tutte le classi $c$. L'\emph{offerta} \`e la
somma delle ore residue dei docenti abilitati a $m$.
Verifica che offerta $\geq$ domanda.
\end{defbox}

Se questo fallisce, l'intera materia \`e in deficit
strutturale: nessuna distribuzione potr\`a coprire tutte le
classi. Soluzione: aggiungere docenti, ridurre il monte ore,
o convocare supplenze.

\subsection{Controllo 3: Hall sampling random-subset}

\begin{defbox}
Si campionano casualmente $K$ sotto-insiemi di lezioni
(tipicamente $K = 200$, di taglie fra 5 e 50) e per ognuno
si verifica direttamente la condizione di Hall.
\end{defbox}

Quando un'istanza ha un \emph{sottoinsieme cattivo} di
lezioni che soddisfano i controlli 1 e 2 ma non Hall, il
campionamento ha alta probabilit\`a di ``inciampare'' su di
esso e segnalarlo. \`E un controllo \emph{euristico} (pu\`o
non beccare casi rari) ma nella pratica $\pi$Tantum lo trova
sufficiente.

\section{Esempio mini-completo}

\begin{esempiobox}
Scuola: 3 classi (1A, 1B, 1C), tutte con 2 ore di
matematica. Due docenti di A027 (Adami e Bruni), Adami con
4 ore residue, Bruni con 3 ore residue. Totale offerta: 7.
Totale domanda: 6. Controllo 2: \textbf{ok}.

Controllo 1: per ogni (classe, mat), 2 ore richieste.
Adami+Bruni hanno entrambi $\geq 2$ ore residue. \textbf{ok}.

Aggiungiamo un vincolo: Bruni non pu\`o insegnare in 1C.
Domanda di 1C-mat: 2 ore. Offerta di docenti compatibili
con 1C-mat: solo Adami con 4 ore residue. \textbf{ok}.
Ma se aggiungiamo un secondo vincolo: Adami satura 3 ore
con altre materie e ne restano solo 1 per le classi 1A/B/C
matematica. Allora la domanda complessiva su Adami+Bruni
per matematica diventa 6 ore, l'offerta 1+3 = 4. Controllo
2 fallisce: si vede subito che servono 2 ore in pi\`u
(rilassare un altro docente, o ridurre il monte).
\end{esempiobox}

\section{L'integrazione nell'UI}

\sidenote{Hall \`e il \emph{primo} controllo che il sistema
esegue. Se fallisce, l'utente vede subito quale gruppo di
classi/materie \`e il problema, prima ancora che la pipeline
parta.}

In $\pi$Tantum il pre-check di Hall \`e esposto in tre punti
dell'interfaccia:
\begin{enumerate}
\item Bottone inline nella card di Phase~A del Workflow.
\item Riga in AdvancedTechniquesCard, con bottone
  ``Esegui'' che lancia tutti e tre i controlli.
\item Sezione del tab Statistiche/Diagnostica, dove il
  risultato \`e mostrato come run asincrono visualizzabile
  in dettaglio.
\end{enumerate}

L'output \`e sempre una lista di ``deficit'' con il
sotto-insieme problematico evidenziato. Esempio:
\begin{lstlisting}
HALL CHECK FAILED
Deficit 1 (controllo 2):
  materia=Matematica, domanda=18, offerta=15
  -- 3 ore mancanti
Deficit 2 (controllo 1):
  classe=4B, materia=Storia: 2 ore richieste,
  ma il docente compatibile (Conti) ne ha solo 1
  residua.
\end{lstlisting}

\section{Limiti e quando non basta}

\begin{avvertenzabox}
\begin{itemize}
\item Il pre-check \`e \emph{necessario} ma non
  \emph{sufficiente}: pu\`o passare anche quando il problema
  \`e infeasibile per ragioni \emph{temporali} (es. lo
  stesso docente serve in due slot incompatibili). Un Hall
  ``ok'' garantisce solo che \emph{le ore totali bastano},
  non che si possano disporre nello spazio settimanale.
\item I controlli 1 e 2 sono esatti; il controllo 3 \`e
  euristico (random sampling). Aumentare $K$ rende il
  controllo pi\`u robusto a discapito del tempo.
\item Hall non considera vincoli logici DSL complessi (es.
  ``il prof X e Y mai stesso giorno''). Quelli vengono
  verificati solo durante l'esecuzione di CP-SAT.
\end{itemize}
\end{avvertenzabox}

\section{Approfondimento: l'algoritmo di Hopcroft-Karp}

\begin{approfondimentobox}
Verificare \emph{esattamente} l'esistenza di un matching che
copre un lato di un grafo bipartito si pu\`o fare in tempo
$O(E\sqrt{V})$ con l'algoritmo di Hopcroft-Karp del
1973~\parencite{hopcroft1973}. \`E quello che usa
internamente OR-Tools per il propagatore globale
\texttt{AllDifferent}: ogni volta che CP-SAT propaga un
vincolo del tipo ``$n$ variabili devono prendere valori
distinti'', costruisce il grafo bipartito \emph{variabili
vs valori} e chiama Hopcroft-Karp per verificare la
fattibilit\`a residua. Se non c'\`e matching $\to$
contraddizione $\to$ backtracking immediato.

In $\pi$Tantum non si usa Hopcroft-Karp esplicitamente nel
pre-check (i tre controlli rapidi sono pi\`u che sufficienti
nei casi pratici), ma l'algoritmo \`e \emph{dentro} CP-SAT
ad ogni propagazione.
\end{approfondimentobox}

\section{Connessioni}

\begin{connessionebox}
\begin{itemize}
\item \textbf{CP-SAT} (cap.~\ref{ch:metodo_cpsat}): la
  propagazione del vincolo \texttt{AllDifferent} in CP-SAT
  usa internamente un calcolo di matching basato sullo
  stesso teorema di Hall.
\item \textbf{Decomposizione spettrale}
  (cap.~\ref{ch:metodo_spettrale}): Hall si esegue \emph{su
  ciascun cluster} oltre che globalmente.
\item \textbf{Workflow}: il pre-check \`e raccomandato come
  \emph{primo} step della pipeline; risparmia minuti di
  CP-SAT su istanze infeasible.
\end{itemize}
\end{connessionebox}

\begin{biblioparvabox}
\begin{itemize}
\item \cite{hall1935}: l'articolo originale di Hall.
  Cinque pagine, una gemma.
\item \cite{hopcroft1973}: l'algoritmo polinomiale per
  matching bipartito, fondante della complessit\`a moderna.
\item \cite{lovasz1986}: \emph{Matching Theory}, libro
  classico per chi vuole approfondire.
\item \cite{papadimitriou1998}, capp.~10--11: sezione
  accessibile sui matching e i flussi.
\end{itemize}
\end{biblioparvabox}
